\scanpage{19}
de Kostant. On a donc des morphismes $f\handwritten{\unsure{\mathrel{\mapstochar\relbar}|\mapsto}}f\handwritten{\circ}i_{\unsure{o|0}}$ et $f\handwritten{\mathrel{\mapstochar\relbar}}f(u_o)$ de $\operatorname{Dyn}_c$ dans $\operatorname{ST}_c$ resp. $\operatorname{Un}_c$. D'autre part, Deligne nous fournit $\operatorname{Un}_c\handwritten{\to\unsure{\underline{\phi}|\underline{\varnothing}}}_c$, et celui-ci peut aussi être défini par la commutativité du diagramme \handwritten{d'esp. homogènes}
\diagram{\begin{tikzcd}[ampersand replacement=\&]
\operatorname{ST}_c\handwritten{\simeq}(\operatorname{CL})_{c'} \handwritten{\arrow[d]} \& \operatorname{Dyn}_c \handwritten{\arrow[l]} \handwritten{\arrow[d]} \\
\handwritten{\underline{\phi}}_c\handwritten{\simeq}\operatorname{Par}_{c'} \& \operatorname{Un}_c \handwritten{\arrow[l]}
\end{tikzcd}}
\margin{\handwritten{\diagram{\begin{tikzcd}[ampersand replacement=\&,column sep=small,row sep=small]
L \arrow[d] \& Z_o \arrow[l] \arrow[d] \\
P \& \unsure{z(e)|Z(e)|\mathfrak{z}(e)} \arrow[l]
\end{tikzcd}}}}
.

Pour qu'il existe un point rationnel de \deleted{$\operatorname{Dyn}_c$, ou ce qui revient au même} $\handwritten{\underline{\phi}}_c\handwritten{\simeq}\operatorname{Par}_{c'}$, il faut (et suffit) qu'il en existe un de $\operatorname{Dyn}_c$ \deleted{i.e. que $G$ soit quasi-déployé}\inserted{\handwritten{(i.e., si $Z_o=e$, que $G$ est quasi-déployé)}}: en effet, $P$ étant choisi, \deleted{\unsure{ceci}} on peut en choisir un \deleted{point rationnel} sous-groupe de Lévi $L$,\deleted{i.e.}\inserted{\handwritten{ou ce qui revient au même}} un $i:\underline{G}_m\handwritten{\to}G$ dans la classe $c$,définissant $P$, et alors les $u=\exp(e)$ qui \deleted{aux} s'adaptent à $i$ forment un ouvert dense de $\handwritten{\unsure{\mathfrak{g}|\mathscr{G}}(2)}$, qui \deleted{est donc non vide} a donc un point rationnel. On voit donc que pour certaines classes de paraboliques (qu'on pourrait appeler les paraboliques de \inserted{Dynkin-}Kostant\unsure{)} l'existence d'un parabolique de cette classe rationnel sur $k$ implique déjà que $G$ est quasi-déployé. (NB Quand $G$ n'est pas une forme intérieure, il faut faire attention cependant de vérifier d'abord que la classe $c'$ de $P$ provient bien d'un $c$ de Dynkin \uline{invariant par} $\handwritten{\Gamma}$, et il n'est pas clair que cette propriété\inserted{de $c$} résulte du fait que $c'$ est \unsure{invariantè|invariante} par $\handwritten{\Gamma}$, faute de savoir si $c\handwritten{\mapsto}c'$ est injectif sur les classes de Dynkin). \handwritten{(}On notera aussi que l'argument donné \deleted{\unsure{p}}marche également sur un anneau semi-local (de car. nulle jusqu'à nouvel ordre).\handwritten{)}
\margin{\handwritten{\unsure{Il|Ce} n'est en effet pas le cas, p.ex. pour $A_2,A_4$ (cela l'est pour $A_3$)}}
Ceci montre l'intérêt, pour des questions arithmétiques, de déterminer exactement les\inserted{classes \handwritten{de}} sous-groupes parab\handwritten{\unsure{-}}oliques de Dynkin-Kostant des divers groupes simples, et de voir si: leur stabilisateur dans $\operatorname{Autext}(G_o)$ est le même que celui des $c$ de Dynkin correspondants, et pas str. plus grand.

On notera cependant que si une filtration du foncteur \deleted{fib} oubli peut être définie par un élément unipotent, la classe de conjugaison de ce der%
% archive sheet number: 18
